\section{Resultados en datos reales}

\subsection{Desempeño agregado}

\begin{table}[H]
\centering
\caption{Rendimiento fuera de muestra en los ocho experimentos reales.}
\label{tab:real_performance}
\begin{tabular}{lccccc}
\toprule
Modelo & RMSE medio & RMSE mediano & $R^2$ medio & Pearson medio & Ganadores \\
\midrule
PINN & \textbf{0,5019} & \textbf{0,5197} & \textbf{0,1525} & \textbf{0,7540} & \textbf{4/8} \\
GLM  & 0,5538 & 0,5387 & 0,0346 & 0,4604 & 2/8 \\
MLP  & 0,5646 & 0,5549 & -0,1418 & 0,7109 & 2/8 \\
\bottomrule
\end{tabular}
\end{table}

La PINN alcanzó el menor RMSE medio y mediano. Ganó cuatro evaluaciones, mientras que GLM y MLP ganaron dos cada uno. No obstante, la ventaja no fue uniforme en todos los escenarios.

\begin{figure}[H]
\centering
\safegraphic[width=\linewidth]{figures/real_models_rmse_by_experiment.png}
\caption{RMSE fuera de muestra de GLM, MLP y PINN en las ocho direcciones reales.}
\label{fig:real_rmse_experiments}
\end{figure}

\subsection{Pruebas estadísticas}

La prueba de Friedman no detectó diferencias globales en RMSE, MAE o $R^2$. Para correlación de Pearson sí se observaron diferencias entre modelos.

\begin{table}[H]
\centering
\caption{Pruebas de Friedman en datos reales.}
\label{tab:real_friedman}
\begin{tabular}{lcc}
\toprule
Métrica & $\chi^2$ & $p$ \\
\midrule
RMSE & 3,000 & 0,2231 \\
MAE & 3,250 & 0,1969 \\
$R^2$ & 3,000 & 0,2231 \\
Pearson $r$ & 9,750 & \textbf{0,0076} \\
\bottomrule
\end{tabular}
\end{table}

Las comparaciones de Pearson mostraron que MLP superó al GLM ($p_{\mathrm{Holm}}=0,0234$) y PINN superó al GLM ($p_{\mathrm{Holm}}=0,0313$). No hubo diferencia significativa entre MLP y PINN ($p_{\mathrm{Holm}}=0,0781$).

\subsection{Brecha de generalización}

La MLP presentó el menor error de entrenamiento, pero la mayor diferencia entre entrenamiento y prueba. La brecha media fue aproximadamente 0,293 puntos porcentuales BOLD para MLP, 0,099 para PINN y 0,063 para GLM.

\begin{figure}[H]
\centering
\safegraphic[width=0.85\linewidth]{figures/real_generalization_gap.png}
\caption{Brecha de generalización definida como RMSE de prueba menos RMSE de entrenamiento.}
\label{fig:real_generalization_gap}
\end{figure}

\subsection{Consistencia de HRF y parámetros}

Las HRF estimadas desde bloques distintos tuvieron correlación media 0,958 y RMSE medio 0,0428 puntos porcentuales BOLD. La forma temporal fue consistente, aunque la amplitud varió en algunos escenarios.

La diferencia simétrica media fue 11,17\% para $\alpha$, 34,97\% para $\epsilon$ y 40,48\% para $\tau$. Por tanto, $\alpha$ fue el parámetro más estable y $\tau$ el menos identificable con un solo bloque.

\subsection{Relación con la repetibilidad de los bloques}

No se observó una asociación concluyente entre correlación de bloques y RMSE. Este análisis utiliza solo cuatro escenarios y debe considerarse exploratorio.

\begin{figure}[H]
\centering
\safegraphic[width=0.9\linewidth]{figures/real_repeatability_vs_rmse.png}
\caption{Relación exploratoria entre repetibilidad de los bloques y error predictivo medio.}
\label{fig:real_repeatability}
\end{figure}
